\section{서론}
\label{sec:introduction}

% 한국어 구조 초안. 논지와 주장 범위만 고정하고, 인용과 문장 수준의
% 다듬기는 영문 본문 작성 단계에서 수행한다.

파운데이션 모델은 전이 가능한 patch feature를 제공하므로, 소수의 normal
support 영상만으로도 few-shot 이상 탐지를 수행할 수 있게 한다
~\cite{oquab2024dinov2,damm2025anomalydino}. 대표적인 구성은 normal support
feature를 static memory에 저장하고, query patch와 가장 가까운 support 사이의
거리를 이상 증거로 해석하는 방식이다~\cite{roth2022patchcore,damm2025anomalydino}.
그러나 이 방식이 강한 표현력을
활용한다는 사실만으로 점수의 신뢰성이 보장되지는 않는다. Normal support와
normal query의 feature가 충분히 비교 가능한 상태로 유지되어야 하기 때문이다.

MVTec AD~2와 같이 학습에 포함되지 않을 수 있는 조명 조건이 test image에
나타나는 환경에서는 이 가정이 쉽게 무너질 수 있다~\cite{hecklerkram2026mvtecad2}.
이를 feature space에서 표현하면
\begin{equation}
  \begin{aligned}
    \widehat P_S^N(\phi) \neq P_Q^N(\phi)
    \quad&\Longrightarrow\quad
    \text{높은 점수}\Leftarrow \\
    &\{\text{분포가 변한 정상},\text{실제 이상}\}.
  \end{aligned}
  \label{eq:score_ambiguity}
\end{equation}
즉, 큰 특징 편차가 반드시 결함에서 비롯되었다고 볼 수 없다. 본 연구의 핵심
전제는 다음과 같다. \emph{강한 표현이 분포 변화 아래에서 신뢰할 수 있는 이상
증거를 자동으로 보장하지는 않는다.}

우리는 이를 \emph{분포 변화 유발 이상 증거 저하(shift-induced anomaly
evidence degradation)}로 정의한다. Few-shot 설정에서는 정상성의 근거가
소수의 support 영상과 그로부터 얻은 patch feature에 제한된다. 이 자체로도
정상 분포의 불완전한 근사를 만들지만, support에 충분히 관측되지 않은 조명,
배치, 재질 외관 및 촬영 조건이 query에 나타나는 support--query distribution
shift 상황에서는 그 불완전성이 더 크게 드러난다. 즉, shift는 새로운 normal
query를 support coverage 밖으로 이동시켜 정상 변화와 실제 결함의 점수 분리를
어렵게 하며, 다음 세 가지 증거 문제를 심화한다.
\begin{equation}
  \begin{aligned}
    \mathcal E_{\mathrm{FS}}
      &:= \text{제한된 few-shot 정상 증거},\\
    \mathcal E_{\mathrm{FS}}
      \xrightarrow{\text{support--query shift}}
      &\{b_{\mathrm{rep}},\ b_{\mathrm{ref}},\ d_{\mathrm{sp}}\},
  \end{aligned}
  \label{eq:evidence_degradation}
\end{equation}
여기서 $b_{\mathrm{rep}}$, $b_{\mathrm{ref}}$, $d_{\mathrm{sp}}$는 각각 표현
편향, 참조 편향 및 공간 증거 저하를 나타낸다.
첫째, \emph{표현 편향}은 동일한 정상 부위의 descriptor가 support와 query에서
일관되지 않게 형성되는 문제이다. ViT feature map에는 positional embedding과
관련된 grid-like artifact 및 high-norm token artifact가 보고되어 왔다
~\cite{yang2024dvt,darcet2024registers}. 따라서 파운데이션 모델의 patch feature는 grid 위치와
주변 문맥의 영향을 받으므로, 배치나 촬영 구도가 달라지면 같은 정상 부위도
서로 먼 특징으로 표현될 수 있다. Few-shot에서는 이러한 변화 양상을 상쇄할
support가 충분하지 않을 수 있으므로, shift가 클수록 정상 패치의 거리가
불필요하게 증가할 수 있다.

둘째, \emph{참조 편향}은 source 정상분포에서 선택된 유한 support가 target
query의 정상성을 판단하는 기준으로 사용될 때 발생하는 coverage mismatch이다.
Few-shot 메모리는 $P_S^N$ 전체가 아니라 몇 장에서 관측된 패치의 경험적
표본이므로, 자주 관측된 정상 영역은 조밀하게 표현하는 반면 드물거나 관측되지
않은 영역은 성기게 표현한다. 실제 memory-bank 계열은 대표 normal patch를
선택해 비교 기준을 구성하며~\cite{roth2022patchcore}, 최근 few-shot 연구도
limited normal prototype의 representativeness를 핵심 문제로 지적한다
~\cite{li2026fastref}. Support와 query가 같은 분포라면 새로운 정상
패치도 메모리가 포괄하는 영역에 놓일 가능성이 상대적으로 높다. 그러나
distribution shift로 $P_Q^N$의 확률 질량이 메모리의 저밀도 또는 미관측 영역으로
이동하면, 의미적으로 정상인 query patch도 가까운 normal support를 찾기 어려워질 수
있다. 이때 최근접 참조 거리와 저밀도 점수는 실제 결함뿐 아니라 reference
coverage의 부족에도 반응하여 정상 패치에 높은 이상 점수를 부여할 수 있다.
따라서 shift가 클수록 참조 자체가 변하는 것이 아니라, 고정된 few-shot
reference와 target 정상분포 사이의 불일치가 커지면서 참조 편향이 더 두드러질
수 있다.

셋째, \emph{공간 증거 저하}는 패치 수준에서 계산된 이상 점수가 현재 query의
정확한 경계와 국소 구조를 충분히 특정하지 못하는 문제이다. 파운데이션 모델의
한 토큰은 일정 영역의 여러 픽셀을 집약하므로, 토큰 내부에서 어느 위치가
이상 반응을 유발했는지에 관한 정보는 제한적이다. 따라서 coarse anomaly
map의 단순 보간은 이미 집약 과정에서 약화된 경계, 얇은 결함 및 국소 대비를
복원할 수 없다. 이러한 공간 해상도 한계는 patch 단위 입력을 사용하는 ViT
구조 자체와 관련되며~\cite{dosovitskiy2021vit}, distribution shift가 없어도
존재하지만, 정상 객체의 배치, 스케일, 촬영 구도 또는 조명이 달라지면 그
영향이 더 두드러질 수 있다. MVTec AD~2가 매우 작은 결함과 학습에 포함되지
않을 수 있는 조명 조건을 함께 평가하도록 설계된 점도 이 문제의 실용적
중요성을 보여준다~\cite{hecklerkram2026mvtecad2}. Shift된 query에서는 객체 경계와 정상 외관의
전이 위치가 support와 다르게 나타날 수 있고, 하나의 토큰이 전경과 배경 또는
서로 다른 국소 외관을 함께 집약할 가능성도 커진다. 이 경우 정상 경계의
응답이 주변으로 번지거나 실제 결함의 응답이 약화되어, coarse map과 실제
결함 영역 사이의 localization 불일치가 커질 수 있다. 즉, shift가 공간 정보를
직접 소실시키는 것이 아니라, 고정된 패치 해상도가 변화된 query 구조를
표현하지 못하는 경우를 늘려 기존 공간 손실의 영향을 심화한다. 현재 query의
RGB 정보는 shift 이후의 경계와 국소 대비를 가장 높은 해상도로 보존하는
관측치이므로, coarse anomaly evidence를 query-native structure에 맞게 정제하는 단서로
사용할 수 있다.

기존 few-shot 이상 탐지 연구는 pretrained feature를 memory matching에
활용하거나~\cite{roth2022patchcore,damm2025anomalydino}, multi-scale feature
aggregation~\cite{jeong2023winclip}, feature reconstruction~\cite{fang2023fastrecon},
또는 prompt learning~\cite{li2024promptad}을 통해 더 판별적인 representation과
matching 성능을 얻는 데 집중해 왔다. 이러한 접근은 정상과 이상을 구분하는 강한 representation을
제공하지만, 정상 support와 query의 표현이 충분히 비교 가능하다는 가정에
의존하는 경우가 많다. 일부 연구가 정합, 증강, 메모리 선택 및 위치 인식
모델링을 통해 shift의 개별 현상을 완화하지만, distribution shift가 정상
표현, few-shot reference coverage 및 공간 localization에 유입시키는 bias를
명시적인 대상으로 함께 다루는 접근은 상대적으로 제한적이다. 본 연구는 강한
representation을 새로 학습하기보다, 동결된 representation에서 이러한
shift-induced bias를 분리해 완화하는 데 초점을 둔다.

이를 위해 동결된 파운데이션 모델 주위에 세 문제와 직접 대응하는 외재적
강건성 스택을 구성한다. 먼저 Patch Positional Debiasing (PPD)은 normal support에서 위치별로 반복되는
descriptor bias를 추정하고 support와 query에서 대칭적으로 제거하여, layout
shift로 인한 representation mismatch를 완화한다. 다음으로 Static
Flow-LatentBank는 보정된 normal support로 Flow를 적합하여 support 분포가
표준화된 latent space를 구성하고, 변환된 support instance를 static bank에
그대로 보존한다. 이를 통해 raw feature space의 불균일한 geometry에 직접
의존하지 않고 latent 1-NN distance로 query를 비교하며, Flow NLL은 weak
one-sided correction으로만 사용한다. 마지막으로 RGB Guide는 fixed patch
resolution에서 손실된 spatial evidence를 보완하기 위해, 현재 query RGB에
남아 있는 boundary와 local contrast를 이용해 coarse anomaly map을
query-native structure에 맞게 정제한다.
모든 support-fitted 구성 요소는 평가 전에 동결되며, 개선 과정은 anomaly label을
사용하거나 메모리를 갱신하지 않는다.

본 연구가 목표로 하는 기여는 다음과 같다.
\begin{itemize}
  \item support--query distribution shift에서 발생하는 높은 점수의 모호성을 분포 변화
        유발 이상 증거 저하로 정식화한다.
  \item support-fitted representation debiasing과 static support-conditioned latent space를 결합한
        외재적 강건성 스택을 제안한다.
  \item 현재 query RGB 영상으로부터 fine-grained spatial information을 학습 없이 복원한다.
  \item PPD, Flow 투영, 잠재 거리, weak NLL, RGB Guide, 자원 조건 및
        주장 범위를 분리하는 대응 실험 프로토콜을 제시한다.
\end{itemize}

본 연구의 범위는 동일 객체 내 배치, 외관, 촬영 조건 및 희소 참조 변화이다.
임의의 카테고리 또는 도메인 변화, 모든 누락 정상 모드의 복원, 기하 정합,
비강체 대응, 완전한 RGB 외관 불변성을 해결한다고 주장하지 않는다. 현재의
공개 MVTec AD~2 결과는 자원 조건이 명시된 shadow evidence이며, 광범위한
강건성과 최종 우월성 주장은 대응된 다중 시드 및 비공개 평가 이후로 유보한다.
